\def\policy{\pi_\theta}
\def\reward{r}
\def\promptDist{\mathcal X}
\def\advantage{\mathcal A}
\newcommand{\prompt}{x}
\newcommand{\response}{y}
\newcommand{\trueanswer}{y^*}
\newcommand{\genacc}{\text{Acc}_\text{gen}(\theta)}
\newcommand{\veracc}{\text{Acc}_\text{ver}(\theta,n)}
\newcommand{\verify}{\text{Verify}}
\newcommand{\verifier}{f}
\newcommand{\rewardver}{R_v}
\newcommand{\rpos}{r_{\text{pos}}}
\newcommand{\rneg}{r_{\text{neg}}}


\section{Preliminaries}


Let $\policy$ denote a language model parameterized by $\theta$. Given a prompt $\prompt$, the model produces a response $\response = (\response^{1}, \response^{2}, \dots)$ auto-regressively. Formally, each token in the response sequence is generated according to the conditional probability:
\begin{equation}
    \response^{k+1} \sim \policy(\cdot \mid \prompt, \response^{\leq k}),
\end{equation}
where we use $\response^{\leq k}$ to refer to the first $k$ tokens generated by the model.

For reasoning-based tasks considered here, the model typically produces responses following a step-by-step ``chain-of-thought'' reasoning approach \citep{wei2022chain}. A verification function $\reward(\response)$, whose dependence on $\prompt$ is omitted for brevity, extracts the model's proposed solution from the generated response and evaluates its correctness against the prompt-specific ground-truth answer:
\begin{equation}\label{eqn.def.reward.binary}
    r(\response) = \begin{dcases}
        1 & \text{if $\response$ is correct,}\\
        0 & \text{if $\response$ is incorrect.}
    \end{dcases}
\end{equation}

Typically, one optimizes the expected \emph{pass rate}, defined as the average accuracy across a distribution of prompts $\promptDist$: 
\begin{equation}\label{eqn.objective}
    J(\theta) = \mathbb{E}_{\prompt \sim \promptDist}\mathbb{E}_{\response \sim \policy(\cdot \mid \prompt)}[\reward(\response)].
\end{equation}

Taking the gradient of the objective function (\ref{eqn.objective}) with respect to $\theta$ and employing a baseline for variance reduction \citep{sutton1998reinforcement} leads to the well-known policy gradient formulation:
\begin{equation}\label{eqn.policy.gradient}
\nabla_{\theta} J(\theta) = \mathbb{E}_{\prompt \sim \promptDist}\mathbb{E}_{\response \sim \policy(\cdot \mid \prompt)}\left[\advantage(\response)\nabla_{\theta}\log \policy(\response \mid \prompt)\right],
\end{equation}
where the advantage function $\advantage(\response)$ is given by:
\begin{equation}\label{eqn.advantage}
\advantage(\response) = \reward(\response) - \mathbb{E}_{\response' \sim \policy(\cdot \mid \prompt)}[\reward(\response')].
\end{equation}
Here $\mathbb{E}_{\response' \sim \policy(\cdot \mid \prompt)}[\reward(\response')]$ is the average pass rate for prompt $\prompt$.
In practice, the policy gradient in ~\eqref{eqn.policy.gradient} is estimated through Monte Carlo samples, yielding the classical REINFORCE algorithm \citep{williams1992simple}. Recent works have modified this base policy gradient formulation to improve its stability, efficiency, and practicality, resulting in advanced methods such as REINFORCE++ \citep{hu2025reinforceefficientrlhfalgorithm}, GRPO \citep{shao2024deepseekmathpushinglimitsmathematical}, PPO \citep{schulman2017proximal}, RLOO \citep{ahmadian-etal-2024-back}, Dr. GRPO~\citep{liu2025understandingr1zeroliketrainingcritical}, GSPO~\citep{zheng2025groupsequencepolicyoptimization}, etc. Implicit to the training method used in our work is the notion of a generation-verification gap \citep{song2025mind}, where generating correct solutions is hard, but verifying them is easy. We present its definition in Appendix \ref{app:gen-ver-gap}.
